% Include in main.tex with: \input{figures/bc_elastic_energy_quasi_elastic}
% Requires in preamble: \usepackage{pgfplots}, \pgfplotsset{compat=1.18}
%                       \usepgfplotslibrary{groupplots}
\begin{figure}[htbp]
\centering
\begin{tikzpicture}
\begin{groupplot}[
    group style={
        group size=2 by 1,
        horizontal sep=1.6cm,
    },
    width=0.46\textwidth,
    height=0.46\textwidth,
    grid=major,
    grid style={gray!25},
    tick align=outside,
    xlabel={Time},
    ylabel={$\displaystyle\int_\Omega \psi_e \, dV$},
    xmin=1.5, xmax=2.0,
    legend style={font=\scriptsize, cells={anchor=west}},
    legend pos=north west,
]

%% --- Left panel: method 1 (old) ---
\nextgroupplot[title={Method \#1}]

\addplot[blue,           thick, solid]
    table[col sep=comma, x=time, y=psie_corr_active_int]{data/bc_sweep/old_method_elastic/reference_fbar_out_1_bc_uniaxial_old.csv};
\addlegendentry{Uniaxial}

\addplot[red,            thick, solid]
    table[col sep=comma, x=time, y=psie_corr_active_int]{data/bc_sweep/old_method_elastic/reference_fbar_out_1_bc_constrained_old.csv};
\addlegendentry{Constrained}

\addplot[blue,           thick, dashed]
    table[col sep=comma, x=time, y=psie_corr_active_int]{data/bc_sweep/old_method_elastic/restart_fbar_out_1_bc_uniaxial_old.csv};

\addplot[red,            thick, dashed]
    table[col sep=comma, x=time, y=psie_corr_active_int]{data/bc_sweep/old_method_elastic/restart_fbar_out_1_bc_constrained_old.csv};

% Vertical marker at restart point
\draw[black!50, dashed] (axis cs:1.6,0) -- (axis cs:1.6,230000);
\node[above, font=\tiny, rotate=90] at (axis cs:1.55,30000) {restart};

%% --- Right panel: method 2 (new) ---
\nextgroupplot[
    title={Method \#2},
    yticklabel pos=right,
]

\addplot[blue,           thick, solid]
    table[col sep=comma, x=time, y=psie_corr_active_int]{data/bc_sweep/new_method_elastic/reference_fbar_out_1_bc_uniaxial_new.csv};

\addplot[red,            thick, solid]
    table[col sep=comma, x=time, y=psie_corr_active_int]{data/bc_sweep/new_method_elastic/reference_fbar_out_1_bc_constrained_new.csv};


\addplot[blue,           thick, dashed]
    table[col sep=comma, x=time, y=psie_corr_active_int]{data/bc_sweep/new_method_elastic/restart_fbar_out_1_bc_uniaxial_new.csv};

\addplot[red,            thick, dashed]
    table[col sep=comma, x=time, y=psie_corr_active_int]{data/bc_sweep/new_method_elastic/restart_fbar_out_1_bc_constrained_new.csv};

\draw[black!50, dashed] (axis cs:1.6,0) -- (axis cs:1.6,230000);
\node[above, font=\tiny, rotate=90] at (axis cs:1.55,30000) {restart};

\end{groupplot}
\end{tikzpicture}
\caption{BC sweep, dynamic elastic case: Integrated elastic energy $\int_\Omega \psi_e \, dV$ vs.\ time near the restart point for method \#1 (left) and method \#2 (right), for the three boundary-condition cases. Solid lines: reference runs. Dashed lines: restart continuations (same color as the corresponding reference, legend in the left panel applies to both panels).}
\label{fig:bc_elastic_energy_quasi_elastic}
\end{figure}
